% SEGMENT OLP-0119-S01
% Part: first-order-logic 
% Chapter: axiomatic-deduction 
% Section: deduction-theorem

% verification of properties of provability needed for maximally 
% consistent sets in the completeness chapter.

\documentclass[../../../include/open-logic-section]{subfiles}

\begin{document}

\iftag{FOL}
      {\olfileid{fol}{axd}{ded}}
      {\olfileid{pl}{axd}{ded}}

\olsection{கழித்தல் தேற்றம்}

% SEGMENT OLP-0119-S02
நாம் கண்டபடி, அடிகோள் அமைப்பில் !!{derivation}s கொடுப்பது கடினமானது;
!!{derivation}s ஐக் கண்டுபிடிப்பதும் கடினமாக இருக்கலாம். நீண்ட !!{formula}s பட்டியல்களை
உண்மையில் எழுதுவதற்குப் பதிலாக, சில எளிய முடிவுகளைப் பயன்படுத்தி அத்தகைய
!!{derivation}s இருப்பதாக வாதிடுவது பொதுவாக எளிது. அத்தகைய மூன்று முடிவுகளை
ஏற்கெனவே நிறுவியுள்ளோம்: $!A \in \Gamma$ என்று அறியும்போது எப்போதும்
$\Gamma \Proves !A$ என்று உறுதிப்படுத்தலாம் என
\olref[ptn]{prop:reflexivity} கூறுகிறது. $\Gamma \Proves !A$ எனில்
$\Gamma \cup \{!B\} \Proves !A$ உம் என்று
\olref[ptn]{prop:monotonicity} கூறுகிறது. மேலும் $\Gamma \Proves !A$ மற்றும்
$!A \Proves !B$ எனில் $\Gamma \Proves !B$ என்று
\olref[ptn]{prop:transitivity} உட்கிடையாகத் தருகிறது. இதோ மற்றோர் எளிய
முடிவு; மோடஸ் போனென்ஸின் ஓர் ``மேல்நிலை'' வடிவம்:

% SEGMENT OLP-0119-S03
\begin{prop}
  \ollabel{prop:mp} $\Gamma \Proves !A$ மற்றும் $\Gamma \Proves !A \lif
  !B$ எனில், $\Gamma \Proves !B$.
\end{prop}

\begin{proof}
  $\{!A, !A \lif !B\} \Proves !B$ என்பது நமக்குத் தெரியும்:
  \begin{derivation}
    1. & $!A$ & Hyp.\\
    2. & $!A \lif !B$ & Hyp.\\
    3. & $!B$ & 1, 2, MP
  \end{derivation}
  \olref[ptn]{prop:transitivity} படி, $\Gamma \Proves !B$.
\end{proof}

% SEGMENT OLP-0119-S04
இந்தச் சூழலில் நாம் பயன்படுத்தவிருக்கும் மிக முக்கியமான முடிவு கழித்தல்
தேற்றம்:

% SEGMENT OLP-0119-S05
\begin{thm}[கழித்தல் தேற்றம்]
\ollabel{thm:deduction-thm} $\Gamma \cup \{!A\} \Proves !B$ எனில்,
அப்போதும் அப்போது மட்டுமே $\Gamma \Proves !A \lif !B$.
\end{thm}

% SEGMENT OLP-0119-S06
\begin{proof}
போதுமான நிபந்தனையை நிறுவும் திசை உடனடியாகப் பின்தொடர்கிறது. $\Gamma \Proves !A \lif !B$ எனில்,
\olref[ptn]{prop:monotonicity} படி $\Gamma \cup \{!A\} \Proves !A \lif
!B$ உம். மேலும் \olref[ptn]{prop:reflexivity} படி $\Gamma \cup \{!A\}
\Proves !A$. எனவே \olref{prop:mp} படி $\Gamma \cup \{!A\} \Proves !B$.

மறுதிசைக்கு, $\Gamma \cup \{!A\}$ இலிருந்து~$!B$ இன் !!{derivation}
நீளத்தின் மீது தொகுத்தறிதல் செய்கிறோம்.

தொகுத்தறிதல் அடிப்படையில், நீளம்~$1$ கொண்ட ஒவ்வொரு !!{derivation} க்கும்
கூற்றை நிறுவுகிறோம். $\Gamma \cup \{!A\}$ இலிருந்து~$!B$ இன் நீளம்~$1$
கொண்ட !!^a{derivation} ஆனது~$!B$ ஐ மட்டுமே கொண்டது; அது சரியானது எனில்
$!B$ ஆனது~$\in \Gamma \cup \{!A\}$ அல்லது ஓர் அடிகோள். $!B \in \Gamma$
அல்லது அது ஓர் அடிகோள் எனில், $\Gamma \Proves !B$. மேலும்
\olref[prp]{ax:lif1} படி $\Gamma \Proves !B \lif (!A \lif !B)$;
\olref{prop:mp} ஆனது $\Gamma \Proves !A \lif !B$ ஐத் தருகிறது. $!B \in
\{ !A\}$ எனில், $\Gamma \Proves !A \lif !B$; ஏனெனில் கடைசி
!!{sentence}~$!A \lif !B$ ஆனது~$!A \lif !A$ க்கு ஒன்றே; இதனை
\olref[pro]{ex:identity} இல் நாம் ஏற்கெனவே !!{derive} செய்துள்ளோம்.

தொகுத்தறிதல் படியில், $\Gamma \cup \{!A\}$ இலிருந்து~$!B$ இன்
!!a{derivation} ஆனது மோடஸ் போனென்ஸால் நியாயப்படுத்தப்பட்ட படி~$!B$ உடன்
முடிகிறது என்க. (அது மோடஸ் போனென்ஸால் நியாயப்படுத்தப்படவில்லை எனில்,
$!B \in \Gamma$, $!B \ident !A$, அல்லது $!B$ ஓர் அடிகோள்; தொகுத்தறிதல்
அடிப்படையில் பயன்படுத்திய அதே வாதம் பொருந்தும்.) அப்போது ஏதோவொரு
!!{formula}~$!C$ க்கு, அந்த !!{derivation}வின் முந்தைய படிகள்~$!C \lif !B$
மற்றும்~$!C$; அதாவது, $\Gamma \cup \{!A\} \Proves !C \lif !B$ மற்றும்
$\Gamma \cup \{!A\} \Proves !C$. அவற்றுக்குரிய !!{derivation}s குறுகியவை;
எனவே தொகுத்தறிதல் கருதுகோள் அவற்றுக்குப் பொருந்தும். ஆகவே பின்வரும்
இரண்டும் கிடைக்கின்றன:
\begin{align*}
  & \Gamma \Proves !A \lif (!C \lif !B); \\
  & \Gamma \Proves !A \lif !C.
\end{align*}
மேலும் \olref[prp]{ax:lif2} படி
\[
\Gamma \Proves (!A \lif (!C \lif !B)) \lif
((!A\lif !C)  \lif (!A \lif !B)),
\]
என்பதும் கிடைக்கிறது. \olref{prop:mp} ஐ இருமுறை பயன்படுத்தினால்,
தேவையானபடி $\Gamma \Proves !A \lif !B$ கிடைக்கும்.
\end{proof}

% SEGMENT OLP-0119-S07
கழித்தல் தேற்றம் செல்லுபடியாகுமாறு துல்லியமாகவே \olref[prp]{ax:lif1},
\olref[prp]{ax:lif2} ஆகியவை தேர்ந்தெடுக்கப்பட்டன என்பதைக் கவனிக்கவும்.

% SEGMENT OLP-0119-S08
!!{derivability} பற்றிய பின்வரும் பயனுள்ள உண்மைகளைப் பயிற்சிகளாக விடுகிறோம்.

% SEGMENT OLP-0119-S09
\begin{prop}
  \ollabel{prop:derivfacts}
\begin{enumerate}
% SOURCE CORRECTION TA-AXD-002: close the outer conditional's parenthesis.
\item $\Proves (!A \lif !B) \lif ((!B \lif !C)
  \lif (!A \lif !C))$; \ollabel{derivfacts:a}
\item $\Gamma \cup \{ \lnot !A\}
  \Proves \lnot !B$ எனில், $\Gamma \cup \{ !B\} \Proves
  !A$ (எதிர்மறைநேர்மாற்று); \ollabel{derivfacts:b}
\item  $\{ !A, \lnot!A\} \Proves
    !B$ (எக்ஸ் ஃபால்சோ குவோட்லிபெட், வெடிப்புக் கொள்கை); \ollabel{derivfacts:c}
\item  $\{ \lnot\lnot!A\} \Proves
  !A$ (இரட்டை மறுப்பு நீக்கம்);\ollabel{derivfacts:d}
\item $\Gamma \Proves \lnot\lnot!A$ எனில், $\Gamma \Proves
  !A$;\ollabel{derivfacts:e}
\end{enumerate}
\end{prop}

% SEGMENT OLP-0119-S10
\tagprob{FOL}
\begin{prob}
  \olref[fol][axd][ded]{prop:derivfacts} ஐ நிறுவுக.
\end{prob}
\tagendprob

\tagprob{notFOL}
\begin{prob}
  \olref[pl][axd][ded]{prop:derivfacts} ஐ நிறுவுக.
\end{prob}
\tagendprob

\end{document}
